
\section{Discrete Hilbert addressing and one-step locality}
\label{app:hilbert_addressing}

The Hilbert curve is a classical space-filling construction mapping an interval continuously onto a square (and, more generally, onto higher-dimensional cubes) \cite{Hilbert1891,Sagan1994}.
For the purposes of this paper, we use only its \emph{finite-resolution} discrete approximation:
a bijection between indices $\{0,\dots,2^{dn}-1\}$ and lattice sites $\{0,\dots,2^n-1\}^d$ that visits each lattice site exactly once.

\paragraph{One-step locality.}
Standard Hilbert orders are constructed so that consecutive indices correspond to lattice neighbors (in $\ell^1$), i.e.\ \eqref{eq:hilbert_locality}.
This property is crucial for interpreting the address map as a locality-preserving scan schedule.

\paragraph{H\"older scaling (continuum limit).}
The continuous Hilbert curve is H\"older continuous with exponent $1/d$ \cite{Sagan1994}, which quantitatively expresses locality preservation under coarse graining.
In the present paper, this scaling is invoked only as a standard mathematical property supporting the finite-resolution interpretation; all operational claims remain at finite $n$.

\paragraph{Implementation (2D).}
For $d=2$, a compact bitwise algorithm maps a Hilbert index to $(x,y)$ coordinates.
We include a pure-Python reference implementation in \path{scripts/exp_hilbert_locality.py}, which also brute-force verifies that $\norm{H_n(t+1)-H_n(t)}_1=1$ for all $t$ and for orders up to $n=8$.
The script emits a small \LaTeX{} row file for Table~\ref{tab:hilbert_locality}.


